\documentclass[11pt,a4paper]{article}

\usepackage[margin=1in]{geometry}
\usepackage{xcolor}
\usepackage{parskip}
\usepackage{enumitem}
\usepackage{hyperref}
\usepackage{framed}
\usepackage{soul}
\usepackage{amsmath,amssymb}
\usepackage{bm}
\usepackage{siunitx}
\usepackage{booktabs}
\usepackage[authoryear,round]{natbib}
\usepackage{xr-hyper}
\usepackage{graphicx}
\usepackage{placeins}    % provides \FloatBarrier to keep figures in section
\usepackage[capitalize,noabbrev]{cleveref}  % \cref for multi-target refs (M1, M4, M5)

\hypersetup{
  colorlinks=true,
  linkcolor=blue!60!black,
  urlcolor=blue!60!black,
  citecolor=blue!60!black,
}

% Pull labels (sections, figures, tables, equations) from the main
% manuscript's .aux file, so that \ref{sec:introduction} etc. resolve
% to the numbering in the compiled manuscript. Requires main.tex to
% have been built first.
% Temporarily disable \bibcite while reading main.aux: otherwise natbib
% sees each citation twice (once from main, once from this letter's bbl)
% and emits "multiply defined" warnings.
\makeatletter
\let\response@orig@bibcite\bibcite
\let\bibcite\@gobbletwo
\externaldocument{../../build/main}
\let\bibcite\response@orig@bibcite
\makeatother

% ---------------------------------------------------------------
% Response-letter environments and macros
% ---------------------------------------------------------------

% Reviewer's comment: left bar, italic, grey
\definecolor{reviewerbar}{gray}{0.55}
\definecolor{reviewerbg}{gray}{0.96}

\newenvironment{rcomment}
  {\begin{leftbar}\color{black!75}\itshape}
  {\end{leftbar}\vspace{0.3em}}

% Author's response: normal style, marked with a bold tag
\newcommand{\response}{\par\noindent\textbf{Response:}\ \ignorespaces}

% Manuscript change pointer (high-level roadmap; the verbatim excerpt
% is shown separately below in the \begin{revtext}...\end{revtext} block)
\newcommand{\changed}[1]{\par\noindent\textbf{Summary of manuscript changes:}\ #1\par}

% Revised manuscript text: light-blue left bar, distinguishes the
% actual added/rewritten manuscript text from the response prose.
% Use after \changed{} to show the new text inline so the reviewer
% does not have to consult the latexdiff PDF.
\definecolor{revtextbar}{rgb}{0.20,0.40,0.70}
\definecolor{revtextbg}{rgb}{0.94,0.96,1.00}
\newenvironment{revtext}
  {\par\vspace{0.2em}\noindent\textit{Revised manuscript text:}\par
   \begingroup\color{black}\begin{leftbar}\small}
  {\end{leftbar}\endgroup\vspace{0.3em}}

% Cross-reference helpers
\newcommand{\manref}[1]{(#1)}                  % e.g., \manref{p.~4, l.~112}
\newcommand{\secref}[1]{Section~#1}
\newcommand{\figref}[1]{Fig.~#1}
\newcommand{\eqnref}[1]{Eq.~#1}

% Number response-letter figures/tables/equations with an R prefix
% (Fig. R1, Fig. R2, ...) so they never collide with manuscript labels.
\renewcommand{\thefigure}{R\arabic{figure}}
\renewcommand{\thetable}{R\arabic{table}}
\renewcommand{\theequation}{R\arabic{equation}}

% Commonly needed symbols in this paper
\newcommand{\Kzero}{\ensuremath{K_0}}
\newcommand{\Zm}{\ensuremath{Z_m}}
\newcommand{\Zmz}{\ensuremath{Z_{m0}}}
\newcommand{\Vs}{\ensuremath{V_s}}
\newcommand{\pmean}{\ensuremath{p^{\prime}}}
\newcommand{\ru}{\ensuremath{r_u}}

% ---------------------------------------------------------------

\title{Response to Reviewers\\
  \large Manuscript COMGE-D-26-01109}
\author{Yusong Han, Ryosuke Uzuoka, Kyohei Ueda}
\date{\today}

\begin{document}

\maketitle

We sincerely thank both reviewers for their thoughtful and constructive
comments. A summary of the major revisions, together with the editor's
cover letter, is provided in the accompanying cover-letter document; the
detailed point-by-point response follows below. In the revised manuscript,
added or modified text is marked using \texttt{latexdiff} against the
submitted version (also provided). All line and figure numbers below refer
to the revised manuscript unless noted otherwise.

\clearpage

%% ===============================================================
%% Reviewer 1
%% ===============================================================

\section*{Response to Reviewer 1}

We thank Reviewer~1 for the detailed and constructive feedback. Our
responses to each point are given below.

%% ---------------- R1.1 ----------------
\subsection*{R1.1 --- Particle size scaling adequacy}

\begin{rcomment}
Lines 112--114: Particle size may influence macroscopic mechanical responses
such as dilatancy, peak friction angle, and liquefaction resistance. Please
clarify the adequacy of this scaling factor (e.g., whether the specimen
size-to-particle size ratio meets the requirement) or provide supporting
references.
\end{rcomment}

\response
We thank the reviewer for raising this important point. As the
reviewer notes, particle size choice can influence macroscopic
responses such as dilatancy, peak friction angle, and liquefaction
resistance. These three behaviours are precisely the targets against
which our calibrated parameter set has been validated (see item~(4)
below). The remainder of this response focuses on the more specific
question of whether the adopted specimen-to-particle size ratio is
adequate. We agree that the radial direction of an HCA specimen
(typically its thinnest dimension) is the direction most susceptible
to specimen-size effects and therefore deserves careful examination.
The present specimen is an annular hollow cylinder with inner
diameter \SI{60}{mm}, outer diameter \SI{100}{mm}, height
\SI{100}{mm}, and a particle-size distribution of
\SIrange{1.0}{3.0}{mm} ($d_{\max}=\SI{3}{mm}$). The corresponding
specimen-to-$d_{\max}$ ratios are 6.7 radially, 33 axially, and 84
in the mid-radius circumferential direction
(Fig.~\ref{fig:rl_ratio}), confirming that the radial ratio is the
binding constraint. We examine its adequacy through the following
independent considerations.

\begin{figure}[htbp]
  \centering
  \includegraphics[width=\linewidth]{figs/annotated/fig_r1_1.pdf}
  \caption{HCA specimen and element detail illustrating the radial
    specimen-to-particle size scaling. The orange wireframe in~(a)
    marks the element region shown enlarged in~(b); the double-headed
    arrow in~(b) shows the radial specimen width $R-r$ against the
    $d_{\max}$ reference scale.}
  \label{fig:rl_ratio}
\end{figure}

\begin{enumerate}[leftmargin=*,label=(\arabic*)]
  \item \textbf{Codified laboratory guideline.} The standardised
    guidelines for triaxial testing both stipulate that the largest
    particle size shall be smaller than one-sixth of the specimen
    diameter: \citet{ASTM_D4767} for consolidated undrained triaxial
    testing on cohesive soils, and \citet{ASTM_D7181} for consolidated
    drained triaxial testing on soils generally (including cohesionless
    sands). The present radial ratio of 6.7 satisfies this 6:1
    threshold; the axial and circumferential ratios satisfy it by a
    wide margin.

  \item \textbf{DEM literature precedent.} Two strands of prior DEM
    work support the present size choice.
    \begin{enumerate}[label=(\alph*), leftmargin=2em, topsep=2pt, itemsep=3pt]
    \item \textit{Size-ratio effect on end-state and cyclic
      behaviour.} \citet{Huang2014} systematically investigated
      the effect of specimen/particle size ratio on DEM triaxial
      samples using a realistic Toyoura-sand grading and concluded
      that ``the ultimate strength of the material at the critical
      state is [\ldots] independent of the sample size and
      boundary conditions''; their smallest sample, at a
      diameter/$d_{\max}$ ratio of 7.2, is comparable to the
      present 6.7. In a more direct cyclic-liquefaction setting,
      \citet{KuhnRenken2014} simulated undrained cyclic shear in a
      DEM assembly with periodic boundaries and minimum cell
      dimension $\sim$12\,$d_{50}$ (comparable to the present
      radial/$d_{50}=10$), and reported that the resulting
      liquefaction resistance curves ($\mathrm{CSR}$ vs $N_L$) were
      nearly indistinguishable when the particle count was doubled.
      Because liquefaction ($\ru=0.95$) is itself an end-state
      condition characterised by vanishing effective stress, the
      quantity reported here ($N_L$) is expected to be weakly
      affected by specimen size in this regime.
    \item \textit{Comparable DEM-HCA configurations.} Within prior
      DEM-HCA studies, \citet{Song2024} simulated a DEM-HCA specimen
      with inner/outer diameters of 60/100~mm (identical to the
      present radial geometry) on upscaled Toyoura sand, using the
      same scaling strategy with an up-scaling factor somewhat
      larger than the 10$\times$ adopted here. \citet{Shi2021} and
      \citet{Liu2021} adopted radial/$d_{\max}$ ratios of 7.5 and
      7.2, respectively, in rigid-wall bonded-particle models of
      sandstone and mudstone under HCA torsional loading, both
      directly comparable to the present 6.7.
    \end{enumerate}

  \item \textbf{HCA-specific geometric argument.} Although the
    radial ratio of 6.7 is the tightest of the three, the primary
    shear stress $\tau_{z\theta}$ and its conjugate shear strain
    $\gamma_{z\theta}$ under torsional loading both develop in the
    axial--circumferential plane, \emph{orthogonal} to the thinnest
    (radial) dimension (Fig.~\ref{fig:rl_shearplane}). The specimen/$d_{\max}$
    ratios in the two directions that carry this primary
    deformation are 33 (axial) and 84 (circumferential), both
    comfortably within the ranges established in~(1) and~(2a).

\begin{figure}[htbp]
  \centering
  \includegraphics[width=\linewidth]{figs/annotated/fig_r1_1_shearplane.pdf}
  \caption{HCA element with cyclic torsional shear $\tau_{z\theta}$
    acting on the $z$-face (top) and $\theta$-face (side). The shear
    plane is the axial--circumferential $(z,\theta)$ plane, orthogonal
    to the (thin) radial direction.}
  \label{fig:rl_shearplane}
\end{figure}

  \item \textbf{Agreement with laboratory data.} The calibrated parameter set
    reproduces both the monotonic undrained response of Toyoura sand
    reported by \citet{Nakata1998} (effective stress path, stress--strain
    relationship, and dilatancy) and the cyclic liquefaction-resistance
    curve of \citet{Ishihara1985}. Agreement across these two independent
    datasets, covering both monotonic and cyclic loading regimes,
    indicates that the adopted particle sizing captures the behaviours
    the model is used to predict.

\end{enumerate}

Taken together, these considerations support the adequacy of the
present size configuration for the quantities reported in this study.

\changed{Section~\ref{subsec:dem_hca} \manref{p.~3, l.~125--131}: two
sentences added immediately after the 10$\times$ scaling clause,
stating the specimen-to-$d_{\max}$ ratios in all three directions,
citing the codified 6:1 threshold and the Huang--Kuhn precedents,
and recording the radial size ratios used by \citet{Shi2021},
\citet{Liu2021}, and \citet{Song2024} as DEM-HCA precedents in the
same range. (These three works also appear in the expanded
Introduction DEM-HCA paragraph under R2.4, grouped by radial
boundary treatment.)}

\begin{revtext}
[\ldots] a scaling approach commonly adopted in DEM simulations of this
sand to reduce particle count while preserving the non-dimensional
grading \citep{Song2024}. \textbf{The resulting specimen-to-$d_{\max}$
ratios are 6.7 (radial), 33 (axial), and 84 (mid-radius
circumferential), satisfying the 6:1 minimum stipulated by
\citet{ASTM_D7181,ASTM_D4767} and consistent with prior DEM studies
reporting weak size-ratio sensitivity of end-state and
cyclic-liquefaction responses at comparable ratios
\citep{Huang2014,KuhnRenken2014}. Prior DEM-HCA studies have adopted
radial size ratios in the same range (7.5 in \citealt{Shi2021}; 7.2
in \citealt{Liu2021}; \citealt{Song2024} used the same 60/100~mm
geometry on upscaled Toyoura sand with a scaling factor somewhat
larger than the 10$\times$ adopted here).} The DEM parameters are
summarized in Table~\ref{tab:dem_parameters}. [\ldots]
\end{revtext}

\FloatBarrier

%% ---------------- R1.2 ----------------
\subsection*{R1.2 --- Void ratio range and ``dense''/``medium-dense'' labels}

\begin{rcomment}
For the case where the authors have defined relative density states of
$D_r \approx 90\%$ and $D_r \approx 75\%$ based on void ratios of 0.736 and
0.742, respectively, the implied range of maximum and minimum void ratios is
approximately 0.732 to 0.772. As a result, I do not think that a void ratio
range of merely 0.04 and the terms ``dense'' and ``medium-dense'' as used in
this manuscript correspond to the physical states observed in the Toyoura
sand that this model purports to represent.
\end{rcomment}

\response
We thank the reviewer for this careful observation. The reviewer's
observation accurately captures a real ambiguity in our labelling:
read as strict point values on Toyoura's tabulated $(e_{\max},
e_{\min})$ scale, ``$D_r \approx 90\%$'' and ``$D_r \approx 75\%$''
would indeed imply an implausibly narrow void-ratio range. The
underlying issue, we believe, is that \emph{relative density is not
a well-defined quantity for a DEM assembly}, and the labels in our
manuscript are behaviour-matched references to laboratory
conditions rather than quantities computed from the DEM void ratio.
The following points explain why we adopted this convention.

\begin{enumerate}[leftmargin=*,label=(\arabic*),topsep=4pt,itemsep=6pt]
\item \textbf{Relative density does not transfer cleanly from laboratory Toyoura to this DEM assembly.}
Two reasons make a direct $D_r$ mapping inappropriate:
\begin{enumerate}[leftmargin=2em,label=(\alph*),topsep=2pt,itemsep=3pt]
  \item \emph{Particle shape mismatch.} The present model uses spherical particles, whereas real Toyoura grains are subrounded-to-subangular; the achievable void-ratio ranges of the two assemblies therefore do not correspond one-to-one.
  \item \emph{No standardised packing-limit procedure for DEM.} The attainable $e_{\max}$ and $e_{\min}$ of a DEM assembly depend on the choice of particle shape and grading, the inter-particle friction used during assembly, the numerical scheme and damping, and the boundary and preparation conditions, none of which are standardised across DEM codes. Relative density therefore lacks the unambiguous reference that it has for laboratory sand.
\end{enumerate}

\item \textbf{The labels are behaviour-matched references to laboratory conditions.}
Rather than computing $D_r$ from the DEM void ratio, we tuned the target void ratio so that the resulting DEM specimen reproduces the macroscopic mechanical response of real Toyoura sand at the corresponding laboratory $D_r$ condition. Specifically:
\begin{enumerate}[leftmargin=2em,label=(\alph*),topsep=2pt,itemsep=3pt]
  \item The ``$D_r \approx 90\%$'' DEM specimen ($e \approx 0.736$)
    simultaneously reproduces (i)~the undrained monotonic stress
    path, stress--strain response, and dilatancy behaviour of
    \citet{Nakata1998}, whose Toyoura specimens were prepared at
    $D_r \approx 92\!-\!93\%$, and (ii)~the cyclic
    liquefaction-resistance curve of \citet{Ishihara1985} for their
    ``very dense'' Toyoura band ($D_r = 88\!-\!95\%$) at
    $K_0 = 1.0$. Both comparisons are shown in
    Fig.~\ref{fig:validation_combined} of the revised manuscript.
  \item The ``$D_r \approx 75\%$'' DEM specimen ($e \approx 0.742$)
    serves as a comparative density anchor for the $K_0$-expansion
    study, chosen so that its macroscopic cyclic response at
    $K_0 = 1.0$ falls within the ``dense'' ($D_r = 75\!-\!82\%$)
    cyclic-resistance data of \citet{Ishihara1985} at representative
    loading amplitude.
\end{enumerate}

\item \textbf{Use of $D_r$-based labels.}
\begin{enumerate}[leftmargin=2em,label=(\alph*),topsep=2pt,itemsep=3pt]
  \item \emph{Convention.} The laboratory data we calibrate against use $D_r$ as the primary specimen label, including \citet{Ishihara1985}'s ``very dense'' ($D_r = 88\!-\!95\%$) and ``dense'' ($D_r = 75\!-\!82\%$) bands. Following the same convention keeps the cross-reference legible to readers of that literature.
  \item \emph{Read as bands, not point values.} Lab $D_r$ is intrinsically a range; \citet[p.~68]{Ishihara1985} explicitly reports ``\emph{a range of $\pm 3.5\%$ around the average}''. Our ``$D_r \approx 90\%$'' and ``$D_r \approx 75\%$'' labels track Ishihara's bands, not exact percentages.
  \item \emph{DEM has no run-to-run variability.} Specimen generation is deterministic and reproducible, so every run of a given preset starts from an identical void ratio. We therefore report a single void ratio per specimen rather than a band, even though the matching laboratory condition spans one.
\end{enumerate}
\end{enumerate}

\changed{
\begin{enumerate}[leftmargin=*,label=(\alph*),topsep=2pt,itemsep=2pt]
  \item Section~\ref{subsec:preparation} \manref{p.~3--4, l.~149--152}: clarified that DEM void ratios cannot be mapped onto Toyoura's textbook $D_r$, owing to the spherical particle shape and the absence of a standardised DEM packing-limit procedure.
  \item Section~\ref{subsec:preparation} \manref{p.~4, l.~152--155}: the labels ``$D_r \approx 90\%$'' and ``$D_r \approx 75\%$'' are now stated as behaviour-matched references to Ishihara's ``very dense'' ($D_r = 88\!-\!95\%$) and ``dense'' ($D_r = 75\!-\!82\%$) density bands.
  \item Sections~\ref{subsec:cyclic_program} \manref{p.~10, l.~278--279} and \ref{subsec:calibration_validation} \manref{p.~11, l.~314} and the caption of Fig.~\ref{fig:validation_combined} \manref{p.~12}: the previous ``dense''/``medium-dense'' verbal labels are replaced by ``very dense''/``dense'', aligning with Ishihara's terminology and removing the wording prone to misreading.
\end{enumerate}
No figure changes in the main manuscript.}

\begin{revtext}
[\ldots] First, the assembly is compressed under zero interparticle
friction to a target void ratio, determined through the iterative
calibration procedure described in
Section~\ref{subsec:calibration_validation}: approximately 0.736 for
the very dense state and 0.742 for the dense state.
\textbf{Because the DEM particles are spherical (in contrast to the
subrounded-to-subangular grains of real Toyoura sand) and because no
standardised procedure exists for determining the packing limits
$(e_{\max}, e_{\min})$ of a DEM assembly, these void ratios are not
converted to a relative density through the textbook formula with
laboratory Toyoura limits. The labels ``$D_r \approx 90\%$'' and
``$D_r \approx 75\%$'' are instead used as behaviour-matched
references to the ``very dense'' ($D_r = 88\!-\!95\%$) and ``dense''
($D_r = 75\!-\!82\%$) density bands of \citet{Ishihara1985}: the DEM
target void ratios are tuned so that the macroscopic monotonic and
cyclic responses reproduce those observed in the laboratory at the
corresponding $D_r$ conditions} (Section~\ref{subsec:calibration_validation}).
[\ldots]
\end{revtext}

\noindent\emph{Terminology relabel} (representative excerpt from the
caption of Fig.~\ref{fig:validation_combined}; analogous edits apply
to Sections~\ref{subsec:numerical_verification} and
\ref{subsec:calibration_validation}):
\begin{revtext}
Calibration and validation of the DEM-HCA model for \textbf{very
dense} Toyoura sand ($D_r \approx 90\%$, $\pmean_0=\SI{100}{kPa}$,
\Kzero{} = 1.0) [\ldots]
\end{revtext}

%% ---------------- R1.3 ----------------
\subsection*{R1.3 --- Missing formulas for CSR and $\Kzero$}

\begin{rcomment}
Several variable symbols are used without a clear computational formula.
Such as CSR, $\Kzero$.
\end{rcomment}

\response
We thank the reviewer for catching these missing definitions. Both
ratios are now defined inline where each symbol first appears.
$\Kzero$ is introduced in Section~\ref{subsec:preparation} at the
opening of the IC--AC consolidation protocol as the ratio of radial
to vertical effective stresses,
$\Kzero = \sigma_r^{\prime}/\sigma_z^{\prime}$; and CSR is introduced
in Section~\ref{subsec:cyclic_program} at the opening of the cyclic
shear program as the amplitude of cyclic shear stress normalised by
the initial mean effective stress,
$\mathrm{CSR} = \tau_{cyc}/p_0^{\prime}$. Both definitions
cross-reference Table~\ref{tab:hca_equations}, which collects the
full set of HCA stress and strain expressions.

\changed{Inline definitions added: $\Kzero$ in
Section~\ref{subsec:preparation} (IC--AC protocol opening
\manref{p.~3, l.~146--147}) and $\mathrm{CSR}$ in
Section~\ref{subsec:cyclic_program} (cyclic-shear program opening
\manref{p.~9, l.~272--273}), with cross-references to
Table~\ref{tab:hca_equations}.}

\begin{revtext}
[\ldots] Specimens are prepared using an isotropic-consolidation to
anisotropic-consolidation (IC--AC) protocol, with the consolidation
anisotropy characterized by the \textbf{lateral-to-vertical effective
stress ratio $\Kzero = \sigma_r^{\prime}/\sigma_z^{\prime}$, where
$\sigma_r^{\prime}$ and $\sigma_z^{\prime}$ are the radial and
vertical effective stresses} (Table~\ref{tab:hca_equations}). [\ldots]
\medskip

[\ldots] Cyclic loading is sinusoidal,
$\tau_{z\theta}(t) = \tau_{cyc}\sin(2\pi f t)$, at frequency
$f = \SI{8}{Hz}$, with the \textbf{cyclic stress ratio defined as
$\mathrm{CSR} = \tau_{cyc}/p_0^{\prime}$ ($p_0^{\prime}$ the initial
mean effective stress after consolidation)}. [\ldots]
\end{revtext}

%% ---------------- R1.4 ----------------
\subsection*{R1.4 --- Consolidated case table}

\begin{rcomment}
The description of the simulation program is summarized in brief statements
without explicitly listing the complete set of parameters. It is recommended
that all simulation cases (including combinations of different densities,
$\Kzero$ values, and CSR levels) be presented collectively in a table
format.
\end{rcomment}

\response
We thank the reviewer for this helpful suggestion. A consolidated
table has been added at the end of
Section~\ref{subsec:cyclic_program} (Table~\ref{tab:case_matrix} of
the revised manuscript, reproduced as Table~\ref{tab:rl_case_matrix}
below). Each row identifies a $(D_r, \Kzero)$ combination; the
third column enumerates the CSR levels tested for that case. The
primary set at $D_r \approx 90\%$ covers
$\Kzero \in \{0.5, 1.0, 2.0\}$ across five CSR levels
($0.200$ to $0.400$ in $0.05$ steps); the expanded set adds two
intermediate $\Kzero$ values ($0.67$ and $1.5$) at $D_r \approx
90\%$ together with the full $D_r \approx 75\%$ strip, all at
CSR $= 0.300$. The resulting $N_L$ values are presented in the
Results section (e.g., Fig.~\ref{fig:liq_resistance_curves}) so
that the methodology table itself remains free of outcome data.

\begin{table}[htbp]
  \centering
  \caption{Cyclic shear simulation program (reproduction of Table~\ref{tab:case_matrix} of the revised manuscript): CSR levels tested at each combination of relative density and consolidation stress ratio. All cases use the displacement-control axial boundary ($dH/dt=0$) matching \citet{Ishihara1985}, except as noted.}
  \label{tab:rl_case_matrix}
  \begin{tabular}{@{}ccl@{}}
    \toprule
    $D_r$ (\%) & $\Kzero$ & CSR levels \\
    \midrule
    90 & 0.50 & 0.200, 0.250, 0.300, 0.350, 0.400 \\
    90 & 0.67 & 0.300 \\
    90 & 1.00 & 0.200, 0.250, 0.300, 0.350, 0.400 \\
    90 & 1.50 & 0.300 \\
    90 & 2.00 & 0.200, 0.250, 0.300, 0.350, 0.400 \\
    \midrule
    75 & 0.50 & 0.300 \\
    75 & 0.67 & 0.300 \\
    75 & 1.00 & 0.300, 0.400$^{\dagger}$ \\
    75 & 1.50 & 0.300 \\
    75 & 2.00 & 0.300 \\
    \bottomrule
  \end{tabular}\\[2pt]
  \begin{minipage}{0.95\linewidth}
    \footnotesize $^{\dagger}$\,Numerical-verification case (Section~\ref{subsec:numerical_verification} of the revised manuscript), run under the constant-total-axial-stress axial boundary.
  \end{minipage}
\end{table}

\changed{New Table~\ref{tab:case_matrix} \manref{p.~10} inserted at
the end of Section~\ref{subsec:cyclic_program}, summarising the full
simulation program with CSR levels enumerated for each
$(D_r, \Kzero)$ case.}

%% ---------------- R1.5 ----------------
\subsection*{R1.5 --- Loading rate / quasi-static condition}

\begin{rcomment}
The loading rate is not specified in the text. In DEM simulations, an
excessively high loading rate may introduce inertial effects, leading to
excessive particle accelerations and distorted contact forces. Please
clarify whether the selected loading rate satisfies quasi-static conditions.
\end{rcomment}

\response
We thank the reviewer for raising this important point. The reviewer
is correct that the loading frequency was not stated explicitly in the
submitted manuscript; this has now been corrected, and the
quasi-static character of the loading has been examined quantitatively
in a new methodology subsection.

\begin{enumerate}[leftmargin=*,label=(\arabic*),topsep=4pt,itemsep=6pt]
\item \textbf{Loading rate.}
The cyclic torsional shear is sinusoidal,
$\tau_{z\theta}(t) = \tau_{cyc}\sin(2\pi f t)$, at frequency
$f=\SI{8}{Hz}$. This value has been added to
Section~\ref{subsec:cyclic_program} of the revised manuscript at the
first definition of CSR.

\item \textbf{Inertial number.}
The standard non-dimensional measure of inertial effects in DEM
\citep{GDRMiDi2004,daCruzEmam2005} is
\begin{equation}
I = \frac{|\dot\gamma_{z\theta}|\,d_{50}}{\sqrt{p^\prime/\rho_s}},
\label{eq:inertial_number_letter}
\end{equation}
with $d_{50}=\SI{2.0}{mm}$ (the median grain diameter of the upscaled
PSD) and $\rho_s=\SI{2650}{kg/m^3}$ (Toyoura particle density).
\citet{GDRMiDi2004} report a strict quasi-static threshold of
$I<10^{-3}$, defined empirically as the regime in which the effective
friction becomes shear-rate independent in steady plane shear; the
companion paper \citet{daCruzEmam2005} subsequently illustrate the
quasi-static regime with the representative value $I=10^{-2}$ in their
Fig.~1(a), giving a second, more permissive reference.

\begin{figure}[htbp]
  \centering
  \includegraphics[width=0.5\linewidth]{figs/fig_r1_5.pdf}
  \caption{Inertial number $I(t)$ for the most demanding case in the
    simulation program ($D_r\approx 75\%$, $\Kzero=1.0$, CSR$=0.400$),
    plotted on log-scale against the normalised cycle ratio
    $N_c/N_L$ ($N_L$ from the $\ru=0.95$ criterion adopted in
    Section~\ref{sec:results} of the revised manuscript). Orange:
    strict quasi-static threshold $I<10^{-3}$ of \citet{GDRMiDi2004}.
    Green: representative quasi-static value $I=10^{-2}$ used by
    \citet{daCruzEmam2005} in their Fig.~1(a). Red dashed: liquefaction
    onset. The same panel forms part~(c) of
    Fig.~\ref{fig:numerical_verification} of the revised manuscript.}
  \label{fig:rl_inertial}
\end{figure}

\item \textbf{Compliance and post-liquefaction excursion.}
We have computed $I(t)$ throughout this run
(Fig.~\ref{fig:rl_inertial}). The pre-liquefaction stage ($N_c/N_L<1$)
lies \emph{entirely} below the representative quasi-static value
$I<10^{-2}$ of \citet{daCruzEmam2005}; the strict threshold $I<10^{-3}$
of \citet{GDRMiDi2004} is satisfied throughout $96.9\%$ of the
pre-liquefaction stage, with the remaining $\sim 3\%$ excursion
confined to the final cycle before liquefaction, where $p^\prime$ has
already decayed by more than an order of magnitude. The loading rate
itself is therefore comfortably quasi-static.

After liquefaction onset, $I$ rises further to peak values of order
$10^{-2}$. We wish to draw the reviewer's attention to the fact that,
unlike monotonic constant-$p^\prime$ shear (the setting in which the
$I<10^{-3}$ threshold was originally calibrated), undrained cyclic
loading is intrinsically a transient $p^\prime$-evolution problem: by
Eq.~(\ref{eq:inertial_number_letter}), $p^\prime\to 0$ at liquefaction
necessarily drives $I$ upward in any test, simulation or laboratory.
The recent DEM-cyclic-liquefaction study by \citet{Wang2025inertial}
is explicit on this point: they identify a three-stage evolution of
$I$---initial stability, rapid transition, and post-liquefaction
stabilisation---and observe that the post-liquefaction increase in
$I$ is a physical signature of the granular fabric reaching a
``kinematic critical state'' rather than an indication that the
simulation has left the dense-flow regime. Our observed peak
post-liquefaction $I\sim 4\times 10^{-2}$ is consistent with this
picture and remains well within the dense-flow regime described by
\citet{daCruzEmam2005}.

\item \textbf{Specimen-level corroboration.}
Independent confirmation that grain-inertia effects are not
contaminating the macroscopic response is provided by the validation
in Section~\ref{subsec:calibration_validation}: the calibrated DEM
specimen reproduces both the monotonic dilatancy and stiffness of
\citet{Nakata1998} and the cyclic liquefaction-resistance curve of
\citet{Ishihara1985} across CSR levels, including the case examined
above. If the loading were dynamically distorted, this combined
agreement would not be expected.
\end{enumerate}

\changed{Section~\ref{subsec:cyclic_program} \manref{p.~9, l.~271--272}: the
loading frequency $f=\SI{8}{Hz}$ and the explicit waveform
$\tau_{z\theta}(t)=\tau_{cyc}\sin(2\pi f t)$ are now stated at the
first appearance of CSR. A new
Section~\ref{subsec:numerical_verification} \manref{p.~10--11, l.~282--311}
(``Numerical verification of the coupled servo'') has been added at
the end of the methodology, with
Fig.~\ref{fig:numerical_verification} \manref{p.~11} reporting the
volume drift, servo tracking errors, and inertial number throughout
a representative post-liquefaction run; the inertial-number panel and
accompanying prose discuss the GDR-MiDi and da Cruz reference values,
the $96.9\%$ / $100\%$ pre-liquefaction compliance, and the physical
origin of the post-liquefaction $I$ excursion. References
\citet{GDRMiDi2004}, \citet{daCruzEmam2005}, and
\citet{Wang2025inertial} have been added to the bibliography.}

\begin{revtext}
[\ldots] The cyclic shear simulation program consists of a verification
case and a parametric \Kzero{} study. \textbf{Cyclic loading is
sinusoidal, $\tau_{z\theta}(t) = \tau_{cyc}\sin(2\pi f t)$, at
frequency $f = \SI{8}{Hz}$}, with the cyclic stress ratio defined as
$\mathrm{CSR} = \tau_{cyc}/p_0^{\prime}$ [\ldots]
\medskip

\textit{(New paragraph in Section~\ref{subsec:numerical_verification},
``Quasi-static condition''):} The inertial number
$I=|\dot\gamma_{z\theta}|\,d_{50}\sqrt{\rho_s/\pmean}$ [\ldots]
\textbf{In the present run, the representative value $I < 10^{-2}$
is satisfied for the entire pre-liquefaction stage ($N_c/N_L < 1$,
$\ru < 0.95$), and the strict threshold $I < 10^{-3}$ is satisfied
throughout $96.9\%$ of that stage.} The remaining $\sim 3\%$
excursion above $10^{-3}$ is confined to the final cycle before
liquefaction [\ldots] \textbf{After liquefaction, $I$ rises further
to peak values of order $10^{-2}$, an unavoidable physical
consequence of the vanishing mean effective stress in the
inertial-number definition that is shared with laboratory undrained
cyclic tests \citep{Wang2025inertial}}; the strict quasi-static
interpretation of $I$ does not apply in this regime in either the
simulation or the laboratory. [\ldots]
\end{revtext}

\FloatBarrier

%% ---------------- R1.6 ----------------
\subsection*{R1.6 --- Rigid wall vs.\ flexible membrane}

\begin{rcomment}
Rigid walls are used to model the inner and outer boundaries of the HCA
specimen, whereas flexible latex membranes are typically used in laboratory
experiments. Rigid boundaries may restrict local deformation and shear band
development, potentially leading to overestimation of stiffness and
liquefaction resistance. Please discuss the limitations of adopting rigid
boundaries and justify their use with relevant references.
\end{rcomment}

\response
We thank the reviewer for raising this point. The concerns the
reviewer lists --- restricted local deformation, suppressed
shear-band development, and potentially overestimated stiffness ---
arise primarily from \emph{drained} HCA testing, where the membrane
locally accommodates the specimen's dilative or contractive response
at developing localisation sites. Under \emph{undrained} loading,
however, the boundary physics is fundamentally different: water
incompressibility imposes a sealed constant-volume condition that
the membrane cannot relax. The choice of rigid cylindrical walls is
required by this physics for the following reasons.

\begin{enumerate}[leftmargin=*,label=(\arabic*),topsep=4pt,itemsep=6pt]
\item \textbf{Undrained physics requires a sealed constant-volume
boundary; the rigid-wall servo is the available DEM construct that
enforces it.}
The laboratory HCA specimen is sealed against drainage, and pore
water is essentially incompressible, so specimen-total volume is
fixed by water conservation regardless of membrane flexibility; the
latex membrane acts only as a pressure interface, not as a
volumetric degree of freedom. The DEM analogue is a global volume
constraint on the specimen boundary, which the present rigid-wall
servo imposes directly at every timestep
(Section~\ref{subsec:cyclic_loading}). Existing flexible-membrane
DEM-HCA formulations \citep{Song2024,Song2024suffusion} were
developed for drained and open-seepage loading, with cell pressure
(or pore-fluid flow) prescribed rather than specimen volume; to our
knowledge, none currently enforces the sealed constant-volume
boundary required here. The rigid-wall-plus-servo configuration is
therefore required by the physics rather than chosen as a
simplification of a readily available flexible-membrane alternative.

\item \textbf{Loading-strain range and validation against laboratory
data.}
The cyclic torsional shear is applied in the $(z,\theta)$ plane,
whereas the rigid-wall constraint acts in the orthogonal radial
direction (cf.\ R1.1). The reviewer's concerns regarding restricted
local deformation and suppressed shear-band development are
characteristic of the regime in which strain localisation develops,
whereas the present liquefaction analysis
terminates at single-amplitude shear strain $\gamma_{SA} = 2.5\%$,
an order of magnitude smaller than the deviatoric strains at which
strain localisation typically develops in laboratory sand under
monotonic shear \citep[e.g.,][]{DesruesViggiani2004}. Stress--strain
agreement with laboratory Toyoura sand has been verified across
this range under both monotonic and cyclic loading: the calibrated
specimen reproduces the undrained monotonic torsional-shear response
of \citet{Nakata1998} up to $\varepsilon_q \approx 5\%$ and the
cyclic CSR--$N_L$ curve of \citet{Ishihara1985} at the
$\gamma_{SA} = 2.5\%$ liquefaction criterion across the full CSR
range (Section~\ref{subsec:calibration_validation}). This
consistency with laboratory data under both loading modes
indicates that any
influence of the rigid-wall boundary on the macroscopic strength
is limited within the strain range that brackets the present
analysis; spatially resolved analysis of localisation features lies
outside the present scope.
\end{enumerate}

\changed{Section~\ref{subsec:dem_hca} \manref{p.~3, l.~111--118}:
a short paragraph has been added at the end of the wall-configuration
description, framing the rigid cylindrical walls as the direct
implementation of the undrained constant-volume boundary via the
servo, and noting that existing flexible-membrane DEM-HCA formulations
were developed for drained (\citealp{Song2024}) and open-seepage
(\citealp{Song2024suffusion}) loading.
\citet{DesruesViggiani2004} added to the bibliography (cited in
this response letter, item~(2) above).}

\begin{revtext}
\textit{(End of wall-configuration paragraph,
Section~\ref{subsec:dem_hca}):} [\ldots] \textbf{Unlike element-level
periodic-cell DEM, which models a representative volume rather than
an apparatus, and unlike recent flexible-membrane DEM-HCA
formulations \citep{Song2024,Song2024suffusion} developed for drained
and open-seepage loading}, the present model uses rigid cylindrical
walls (inner and outer cylinders, top and bottom planes) with six
torsional blades to apply shear.
\medskip

\textbf{The undrained condition is reproduced by enforcing constant
specimen volume directly through the rigid-wall servo at every
timestep (Section~\ref{subsec:cyclic_loading}). This is the same
kinematic constraint that, in the laboratory, follows from the
essentially incompressible pore fluid: the latex membrane transmits
cell pressure to the grain skeleton but cannot change the specimen
volume. Rigid walls are accordingly an effective DEM implementation
of the undrained boundary.}
\end{revtext}

%% ---------------- R1.7 ----------------
\subsection*{R1.7 --- Citation for CSL in Fig.~8(c)}

\begin{rcomment}
The CSL in Fig.~8(c) is of unknown origin; whether it comes from your own
data or from someone else's literature, a clear citation is required.
\end{rcomment}

\response
We thank the reviewer for raising this point. We agree that the
dashed line in Fig.~\ref{fig:macro_response_c}(c) needs an explicit
source.

On re-examining the slope used in the submitted version, we found
it had been inferred from the endpoint of our calibration
monotonic-undrained-torsional-shear simulation at
$\varepsilon_q\approx 5\%$. This endpoint is not guaranteed to lie
on the critical-state plateau---where $q/p^{\prime}$ has fully ceased
to evolve with strain---so the slope inferred at that point cannot
rigorously be presented as the critical-state slope. We have
therefore replaced it with the literature value of
\citet{VerdugoIshihara1996} for Toyoura sand. The figure caption
now identifies the dashed line as the critical-state line at slope
$M_{cs}\approx 0.892$, and the main text immediately following
Fig.~\ref{fig:macro_response_c} states the citation together with
the Mohr--Coulomb conversion to the present pure-shear loading.

The canonical drained-triaxial-compression dataset of
\citet{VerdugoIshihara1996} for Toyoura sand reports a steady-state
friction angle $\varphi_{cs}\approx 31^{\circ}$. Because their
measurement is at intermediate principal stress ratio $b=0$ whereas
the HCA torsional boundary conditions of the present work impose
pure-shear loading by construction ($b=0.5$), the slope of the CSL
in $p^{\prime}$--$q$ space is converted via the Mohr--Coulomb
relation
\[
  M_{cs} = \sqrt{3}\,\sin\varphi_{cs}
  \approx \sqrt{3}\,\sin(31^{\circ})
  \approx 0.892,
\]
which is the value at which the dashed line in
Fig.~\ref{fig:macro_response_c}(c) is drawn.

\changed{Caption of Fig.~\ref{fig:macro_response_c}(c)
\manref{p.~14}: the dashed line is now identified as the
critical-state line at slope $M_{cs}\approx 0.892$ (with the
attribution to \citet{VerdugoIshihara1996} and the Mohr--Coulomb
conversion to $b=0.5$ stated in the main text immediately
following the figure, Section~\ref{subsec:macro_response}
\manref{p.~13, l.~391--393}).}

\begin{revtext}
\textit{Updated caption of Fig.~\ref{fig:macro_response_c}(c):}
Deviatoric stress vs.\ $p^{\prime}$ for three \Kzero{} states;
\textbf{the dashed line is the critical-state line at slope
$M_{cs}\approx 0.892$ (see text)}.

\medskip

\textit{New main-text sentence accompanying
Fig.~\ref{fig:macro_response_c}(c) in
Section~\ref{subsec:macro_response}:}
The dashed line shown in panel~(c) for reference is the
\textbf{critical-state line at $M_{cs}\approx 0.892$, obtained
from the steady-state friction angle
$\varphi_{cs}\approx 31^{\circ}$ reported by
\citet{VerdugoIshihara1996} for Toyoura sand via
$M_{cs}=\sqrt{3}\sin\varphi_{cs}$ at $b=0.5$.}
\end{revtext}

%% ---------------- R1.8 ----------------
\subsection*{R1.8 --- ``Enhancement'' of $\Kzero = 1.0$ in the coordination-number plane (Fig.~12)}

\begin{rcomment}
In most cases, $N_L$ decreases monotonically with the increase of $\Kzero$,
as shown in Figures 8 and 10; however, the situation changes in Figure 12.
Theoretically, the evolution of the coordination number at the microscopic
level corresponds to the evolution of the macroscopic parameters such as the
excess pore water pressure ratio, deviatoric stress, and cumulative shear
work. The enhancement of $\Kzero = 1.0$ in the coordination number plane
should be explained more clearly.
\end{rcomment}

\response
We thank the reviewer for this observation; it is right to flag this
apparent paradox. We confirm directly from the
simulation data that
$Z_{m0}|_{\Kzero=1.0}=4.978 > Z_{m0}|_{\Kzero=0.5}=4.894 \approx
Z_{m0}|_{\Kzero=2.0}=4.885$ at $D_r = 90\%$ — i.e.\ the
isotropically consolidated specimen does carry the highest initial
coordination number — while $N_L$ in the same dataset (inset of
Fig.~\ref{fig:zm_evolution}) decreases monotonically from
\Kzero{} = 0.5 through 1.0 to 2.0. The two rankings are therefore genuinely
different, not a plotting artefact, and the apparent paradox is in
fact the central microscopic finding of the paper rather than an
inconsistency: \Zm{} alone is not sufficient to predict liquefaction
resistance under torsional shear.

The physical reason is that \Zm{} is a scalar count over all contact
orientations, so it favours the unbiased packing produced by
isotropic consolidation; under \Kzero{} = 1.0 there is no preferred
direction along which contact loss is concentrated during
consolidation, allowing slightly more contacts to be retained than
under \Kzero{} $\neq 1$. Resistance to torsional shear $\tau_{z\theta}$,
however, depends on which contacts carry load in the $(z,\theta)$
plane: it is the \emph{directional partitioning} of the contact
network that determines $N_L$. This is examined quantitatively in
Section~\ref{subsec:micro_fabric_results}, where both the contact-
density (count-weighted) and normal-force (force-weighted,
Fig.~\ref{fig:contact_force_sequence}) analyses consistently identify
the fabric anisotropy as the relevant microscopic descriptor: $\alpha_0$ correlates positively with $N_L$
within each density group (Fig.~\ref{fig:fabric_liq_pair}b), and the
torsion-resisting fractions $1-\Phi_{rr}$ and $1-\Phi^f_{rr}$ both
decrease monotonically with \Kzero{}
(Fig.~\ref{fig:fabric_mechanism}, Table~\ref{tab:fabric_diagonal}).

To prevent the same confusion at first reading, we have added a brief
forward-pointing sentence to the paragraph discussing
Fig.~\ref{fig:zm_evolution}
(Section~\ref{subsec:micro_fabric_results}) noting that the highest
$Z_{m0}$ does not coincide with the highest $N_L$, with a one-clause
hint (count vs.\ directional partitioning) and a cross-reference to
Fig.~\ref{fig:fabric_liq_pair} and Fig.~\ref{fig:fabric_mechanism}
where the resolution is given. The full quantitative argument is
deliberately not duplicated at Fig.~\ref{fig:zm_evolution}, to
preserve the build-up from \Zm{} (insufficient) $\to$ $\alpha$
(necessary) $\to$ $\Phi^f$ (dominant) that frames the microscopic
discussion.

\changed{Forward-pointing sentence added at
Fig.~\ref{fig:zm_evolution} discussion in
Section~\ref{subsec:micro_fabric_results}
\manref{p.~17, l.~462--466}; full resolution remains in the same
section at Fig.~\ref{fig:fabric_liq_pair} \manref{p.~22} and
Fig.~\ref{fig:fabric_mechanism} \manref{p.~23} (the latter expanded
under R2.7).}

\begin{revtext}
[\ldots] the isotropically consolidated specimen (\Kzero{} = 1.0)
exhibits the largest initial value ($Z_{m0} \approx 4.98$), whereas
the anisotropically consolidated specimens (\Kzero{} = 0.5 and 2.0)
have slightly smaller and nearly identical values
($Z_{m0} \approx 4.89$ and $4.88$, respectively).
\textbf{Yet \Zm{} alone cannot rank the three: the \Kzero{} = 1.0
specimen carries the largest $Z_{m0}$ but is not the most resistant
(Fig.~\ref{fig:liq_resistance_curves}), and the \Kzero{} = 0.5 and
2.0 specimens share nearly identical $Z_{m0}$ yet have markedly
different $N_L$ (inset of Fig.~\ref{fig:zm_evolution}). The reason
is that \Zm{} counts contacts isotropically, whereas resistance to
torsional shear depends additionally on the directional partitioning
of contacts and forces, examined quantitatively in the paragraphs
that follow.} [\ldots]
\end{revtext}

\clearpage

%% ===============================================================
%% Reviewer 2
%% ===============================================================

\section*{Response to Reviewer 2}

We thank Reviewer~2 for the thorough and constructive comments, and
particularly for pointing us to relevant prior work on DEM-HCA modelling.

%% ---------------- R2.1 ----------------
\subsection*{R2.1 --- Particle--wall friction coefficient set to zero (Table 1)}

\begin{rcomment}
Table 1: Why is the friction coefficient between particles and walls zero?
\end{rcomment}

\response
We thank the reviewer for raising this point. The zero
particle--wall friction coefficient applies uniformly to all
wall-type elements in the PFC model: the inner and outer
cylinders, the top and bottom platens, and the six torsional blades
(which are also wall elements in PFC). The physical justification
differs by wall class, and we organise the response accordingly.

\begin{enumerate}[leftmargin=*,label=(\arabic*),topsep=4pt,itemsep=6pt]
\item \textbf{Cylindrical walls --- principal-plane preservation.}
The HCA element-test reduction (Appendix~A,
Table~\ref{tab:hca_equations}) treats the radial-normal face (the
inner and outer cylinders) as a principal stress plane, which
requires zero shear traction on it. The laboratory satisfies this
automatically: under quasi-static loading the flexible membrane is
in mechanical equilibrium with the surrounding pressurised water,
and its sand-side tangential tractions integrate to zero at the
element scale. DEM walls are massless and carry no such
self-equilibrium; any $\mu_{pw}>0$ would leave a residual shear
resultant on the radial face (Fig.~\ref{fig:rl_forces}), violating
the principal-plane assumption on which the reduction rests.
$\mu_{pw}=0$ on the cylinders is therefore \emph{required}, not
merely convenient.

\begin{figure}[htbp]
  \centering
  \includegraphics[width=\linewidth]{figs/annotated/fig_r2_1.pdf}
  \caption{Particle--wall contact-force decomposition at the outer
    cylindrical wall. (a) marks the element enlarged in (b), where
    arrows show one representative contact: red $f_{pw_n}^{o}$ is the
    normal component (the only one present with $\mu_{pw}=0$); the
    ghost blue $f_{pw_t}^{o}$ is the tangential component a non-zero
    $\mu_{pw}$ would introduce. The macroscopic $\tau_t$ is the
    area-average of all $f_{pw_t}^{o}$, so $\mu_{pw}=0$ enforces
    $\tau_t=0$ exactly. Notation: Appendix~A,
    Table~\ref{tab:hca_equations}.}
  \label{fig:rl_forces}
\end{figure}

\item \textbf{End platens and blades --- geometric torque
transmission.}
The axial-normal face is not a principal plane: $\tau_{z\theta}$
acts on it, and the HCA must transmit this shear from the end
platens into the specimen. The laboratory rotates the end platens,
combining sand--cap friction with platen-mounted blades. Mimicking
this friction in DEM with $\mu_{pw}>0$ introduces a servo
instability: the torque transmitted through the platens then depends
simultaneously on the normal and tangential contact forces, coupled
by the Coulomb bound $|f_{pw_t}^{t}|\le\mu_{pw}\,f_{pw_n}^{t}$
(Fig.~\ref{fig:rl_platen_forces}). The servo must track both
contributions at once, and this coupling becomes pathological near
liquefaction, when $f_{pw_n}^{t}$ decays with $\sigma_z^{\prime}$
and contacts detach.

\begin{figure}[htbp]
  \centering
  \includegraphics[width=\linewidth]{figs/annotated/fig_r2_1_platen.pdf}
  \caption{Particle--wall contact-force decomposition at the top
    platen. (a) marks the top element enlarged in (b), where arrows
    show one representative contact (clear of the blade at
    $\theta=0$): red $f_{pw_n}^{t}$ is the normal component; the
    ghost blue $f_{pw_t}^{t}$ is the tangential component a non-zero
    $\mu_{pw}$ would introduce, bounded by
    $|f_{pw_t}^{t}|\le\mu_{pw}\,f_{pw_n}^{t}$. With $\mu_{pw}=0$ each
    $f_{pw_t}^{t}=0$, so torque is transmitted through normal
    contacts on the blade faces rather than through Coulomb-limited
    platen friction (whose budget vanishes with $\sigma_z^{\prime}$
    near liquefaction). Notation: Appendix~A,
    Table~\ref{tab:hca_equations}.}
  \label{fig:rl_platen_forces}
\end{figure}

The present model sidesteps this by transmitting torque
\emph{geometrically} through the six blades rather than through
platen friction. Each blade is a thin radial slab whose face normal
points circumferentially, so a normal particle--blade contact
delivers a purely circumferential force (the torque-transmitting
direction) with no tangential friction required. Setting
$\mu_{pw}=0$ on both platens and blades therefore cleanly separates
the two budgets: pressure is carried by normal contacts on the
platens and cylinders, and torque is carried by normal contacts on
the blade faces.
\end{enumerate}

\changed{Section~\ref{subsec:dem_hca} \manref{p.~3, l.~138--142}:
added a clause after the torsional-blade paragraph stating that the
cylindrical walls and end platens are assigned $\mu_{pw}=0$ so that
they act as pure pressure boundaries, with torque transmission
provided geometrically by the blade faces, and that on the cylindrical
walls in particular this is required by the HCA element-test reduction
(Appendix~\ref{sec:appendix_hca}, Table~\ref{tab:hca_equations}), which
treats the radial-normal face as a principal-stress plane. The full
Coulomb-limit argument for the platens (Fig.~\ref{fig:rl_platen_forces})
is kept in this letter rather than added to the manuscript, since it
concerns a platen-design rationale orthogonal to the paper's focus on
$K_0$ effects and fabric evolution.}

\begin{revtext}
[\ldots] these independent wall-based blades avoid interference with
the radial expansion or contraction of the cylinders, thereby
preserving accurate radial stress measurement. \textbf{The
cylindrical walls and end platens are assigned zero friction
($\mu_{pw}=0$, Table~\ref{tab:dem_parameters}) so that they act
purely as servo-controlled pressure boundaries; torque is
transmitted geometrically through normal contact forces on the
blade faces and does not rely on a wall-friction coefficient. On
the cylindrical walls in particular, $\mu_{pw}=0$ is also required
for the HCA element-test reduction (Appendix~\ref{sec:appendix_hca}),
which treats the radial-normal face as a principal-stress plane:
with $\mu_{pw}>0$, torsion would generate a residual circumferential
particle--wall traction on this face and thereby invalidate the
principal-direction assumption.}
\end{revtext}

\FloatBarrier

%% ---------------- R2.2 ----------------
\subsection*{R2.2 --- Size of torsional blades}

\begin{rcomment}
What is the size of torsional blades?
\end{rcomment}

\response
We thank the reviewer for asking. The blade height is
\SI{15}{mm} (Fig.~\ref{fig:rl_blade_height}), and
the blades extend radially across the full annulus between the inner
and outer cylinders. Six blades are inserted from the top and bottom
plates (three each), uniformly spaced in the circumferential
direction; their arrangement is visible in
Fig.~\ref{fig:specimen_generation}.

\begin{figure}[htbp]
  \centering
  \includegraphics[width=0.55\linewidth]{figs/annotated/fig_r2_2.pdf}
  \caption{Blade height $h_b$ annotated on the DEM-HCA
    specimen during air-pluviation (cf.\
    Fig.~\ref{fig:specimen_generation}(a)); $h_b=\SI{15}{mm}$ in
    the present model.}
  \label{fig:rl_blade_height}
\end{figure}

The blades of the laboratory HCA apparatus are approximately
\SI{1.5}{mm}, which is \SI{15}{} times the minimum particle size of
Toyoura sand. In the present DEM model, the blade height is
correspondingly set to \SI{15}{} times the minimum particle size of
the DEM grading ($d_{\min}=\SI{1.0}{mm}$, see the
\SIrange{1.0}{3.0}{mm} grading stated in
Section~\ref{subsec:dem_hca}), giving \SI{15}{mm}. This preserves
the blade-to-grain scaling of the laboratory apparatus.

\changed{Section~\ref{subsec:dem_hca} \manref{p.~3, l.~133--135}:
the blade-insertion sentence now states the blade height (\SI{15}{mm})
explicitly, followed by a short sentence giving the blade layout
(three blades at each end of the specimen, extending radially across
the annulus and uniformly spaced around the specimen axis).
Figure~\ref{fig:specimen_generation}(a) \manref{p.~4} is annotated
with a dimension arrow labelling $h_b$ on a visible
top-end blade (the value $h_b=\SI{15}{mm}$ is stated in the main
text).  The blade-to-grain-size rationale
is kept in this letter rather than added to the manuscript: the
methodology section reports the blade height as a model parameter,
but the paper does not otherwise discuss whether the
blade-to-$d_{\min}$ ratio is appropriate for DEM-HCA modelling,
which is orthogonal to the paper's focus on $K_0$ effects and
fabric evolution.}

\begin{revtext}
[\ldots] Gravity is then removed to ensure a uniform stress state,
and \textbf{six torsional blades, each 15~mm tall, are inserted into
the specimen} (Fig.~\ref{fig:specimen_generation}b). \textbf{Three
blades at each end of the specimen extend radially across the
annulus between the inner and outer cylinders, uniformly spaced
around the specimen axis.} [\ldots]
\end{revtext}

%% ---------------- R2.3 ----------------
\subsection*{R2.3 --- Difference vs.\ servo mechanism in Han (2024)}

\begin{rcomment}
What is the difference between the servo mechanism used in this paper and
that in Han (2024)?
\end{rcomment}

\response
We thank the reviewer for prompting us to make this comparison
explicit. \citet{Han2024jgs} held the specimen height fixed
($dH/dt=0$),
whereas the present servo treats $dH/dt$ as an independent unknown
solved at each timestep. Three differences follow:
\begin{enumerate}[leftmargin=*,label=(\alph*),topsep=2pt,itemsep=3pt]
  \item \textbf{Kinematic scope.} With $dH/dt=0$, the
    constant-volume constraint eliminates $dR/dt$ in favour of
    $dr/dt$, leaving a single radial unknown. The present
    formulation retains $dH/dt$, so two unknowns ($dr/dt$ and
    $dH/dt$) remain after the volume constraint is applied.
  \item \textbf{Cross-coupling.} A single-unknown system is
    necessarily scalar, with no off-diagonal terms to resolve. With both $dr/dt$ and $dH/dt$ active, the two rates
    jointly affect both the radial and the axial stresses, so the
    servo coefficients form a $2\times 2$ coupled stiffness matrix
    $\hat{\mathbf{K}}$. We derive it analytically and invert it in
    closed form at every timestep, with the off-diagonal entries
    $S_{rz}$, $S_{zr}$ representing the mutual influence of $dr/dt$
    and $dH/dt$.
  \item \textbf{Boundary-condition flexibility.} \citet{Han2024jgs}
    held the radial pressure-difference target at
    $p_{dif,r}^{tar}=0$, effectively assuming equal inner and outer
    cell pressures ($p_o=p_i$). The laboratory HCA, however,
    controls $p_o$, $p_i$, and $p_z$ through independent channels
    (Fig.~\ref{fig:rl_hca_schematic}), so $p_{dif,r}^{tar}$ is a free
    input that can in general be an arbitrary function of time. The
    present servo treats $p_{dif,r}^{tar}$ as such a general input,
    recovering the Han~(2024) condition as the special case
    $p_{dif,r}^{tar}\equiv 0$.
\end{enumerate}

\begin{figure}[htbp]
  \centering
  \includegraphics[width=\linewidth]{figs/annotated/fig_r2_3.pdf}
  \caption{Laboratory HCA (left) and DEM constant-volume
    idealisation (right), reproduced from
    Fig.~\ref{fig:hca_schematic} of the manuscript. The outer cell
    pressure $p_o$, inner cell pressure $p_i$, and axial pressure
    $p_z$ act through independent channels, so they may be
    prescribed as general functions of time; the pressure-difference
    targets $p_{dif,r}^{tar}(t)=p_o(t)-p_i(t)$ and
    $p_{dif,z}^{tar}(t)=p_z(t)-\sigma_r(t)$ inherit this
    time-variability.}
  \label{fig:rl_hca_schematic}
\end{figure}
Re-imposing $dH/dt=0$ and $p_{dif,r}^{tar}\equiv 0$ collapses
$\hat{\mathbf{K}}$ to its radial-only scalar form, so the present
formulation reduces exactly to \citet{Han2024jgs} in that limit
(stated after Eqs.~(\ref{eq:servo_dr})--(\ref{eq:servo_dH}) in
Section~\ref{subsec:cyclic_loading}).

\changed{Section~\ref{subsec:cyclic_loading}: the laboratory HCA
paragraph (around Fig.~\ref{fig:hca_schematic} \manref{p.~6}) now
states explicitly that $p_o$, $p_i$, and $p_z$ are controlled through
independent channels and can be prescribed as general functions of
time; the contribution paragraph
\manref{p.~9, l.~259--270} correspondingly describes
$p_{dif,r}^{tar}(t)$ and $p_{dif,z}^{tar}(t)$ as general
time-dependent inputs, and the comparison with \citet{Han2024jgs}
is updated to note their fixed $p_{dif,r}^{tar}=0$ (i.e.,
$p_o=p_i$) alongside the fixed $dH/dt=0$.
Eq.~(\ref{eq:epwp_definition}) \manref{p.~9, l.~254--258} is
generalized to $u = \sigma_r - \sigma_r^{\prime}$ (mean radial total
minus mean radial effective stress) so that the excess pore water
pressure definition remains valid under a time-varying
$p_{dif,r}^{tar}$; it reduces to the previous
$u=\sigma_{r0}^{\prime}-\sigma_r^{\prime}$ for the constant-cell-
pressure cases studied here. For transparency, the cyclic-shear
program description \manref{p.~10, l.~274--278} now states explicitly
that the cell pressures are held at $p_o=p_i=\sigma_{r0}$
(i.e., $p_{dif,r}^{tar}=0$) and the axial boundary is imposed in
displacement-control mode ($dH/dt=0$) for the cyclic $K_0$ suite
that forms the core of this paper. The monotonic validation
runs against \citet{Nakata1998} and a supplemental cyclic
validation set (added in response to R2.6 to demonstrate exact
volume conservation under a free axial boundary) both use the
constant-total-axial-stress formulation, with principal-angle-
locked $p_{dif,z}^{tar}(t)$ in the monotonic case. The general
$p_{dif,r}^{tar}(t)$ input remains available in the servo
code. We thank the
reviewer for prompting these clarifications, which make the
generality of the boundary-condition definition explicit in the
manuscript.}

\begin{revtext}
\textit{(Laboratory-HCA paragraph,
Section~\ref{subsec:cyclic_loading}):} In a laboratory HCA
[\ldots] the specimen is enclosed between an outer chamber (pressure
$p_o$) and an inner chamber (pressure $p_i$), \textbf{which are
controlled through independent channels and can therefore be
prescribed as general functions of time}. An axial load is applied
through the loading piston, producing \textbf{a vertical pressure
$p_z$ that is likewise independently prescribed} [\ldots]
\medskip

\textit{(Generalised pore-pressure definition,
Eq.~(\ref{eq:epwp_definition})):}
\textbf{$u = \sigma_r - \sigma_r^{\prime}$}, where $\sigma_r$ is the
mean radial total stress prescribed by the laboratory boundary
conditions and $\sigma_r^{\prime}$ is the current mean radial
effective stress evaluated from the DEM contact forces. \textbf{For
the constant-cell-pressure cases studied here, $\sigma_r$ is
time-invariant and equal to $\sigma_{r0}^{\prime}$, so
\cref{eq:epwp_definition} reduces to
$u = \sigma_{r0}^{\prime} - \sigma_r^{\prime}$.} [\ldots]
\medskip

\textit{(Comparison-with-Han2024 clause, contribution paragraph):}
[\ldots] \textbf{\citet{Han2024jgs} solved only the radial and
torsional conditions while keeping the specimen height fixed
($dH/dt=0$) and the radial pressure-difference target held at
$p_{dif,r}^{tar}=0$ (i.e., $p_o=p_i$)} [\ldots]
\end{revtext}

%% ---------------- R2.4 ----------------
\subsection*{R2.4 --- DEM-HCA literature review and the $H=D$ height choice}

\begin{rcomment}
There should be more context to introduce the research in DEM HCA models in
Part Introduction. \dots\ [boundary-condition types in Li~2014, Li~2017,
Song~2024a, Song~2024b; height $H \approx 2D$ in Shi~2021]. However, the
height equals to the outer diameter in this paper. Are there any
requirements or guidelines for the size of HCA?
\end{rcomment}

\response
We thank the reviewer for pointing us to this body of prior DEM-HCA
work. Our response has two parts: an expanded introduction of prior
DEM-HCA research, and a justification of the $H=D$ aspect ratio.

\begin{enumerate}[leftmargin=*,label=(\arabic*),topsep=4pt,itemsep=6pt]
\item \textbf{Prior DEM-HCA research and boundary taxonomy.}
The submitted manuscript has been revised to cite the DEM-HCA
studies identified by the reviewer, grouped by the radial
boundary-condition treatment that distinguishes them. Within the
\emph{stacked-wall} family (mimicking membrane compliance through
stacks of small wall segments), \citet{LiGuoZhang2014} applied
cyclic torsional loading under drained conditions,
\citet{LiZhangGutierrez2015} examined the HCA torsional shear test
itself numerically, and \citet{LiChenGutierrez2017} extended this
approach to drained loading under varying intermediate principal
stress ($b$) and principal stress direction ($\alpha$). Within the
\emph{rigid-wall} family, \citet{Shi2021} and \citet{Liu2021}
combined rigid cylindrical walls with bonded-particle assemblies
to analyse rock-like HCA behaviour. The \emph{flexible
bonded-sphere membrane} family, most recently developed by
\citet{Song2024} for drained sand and by \citet{Song2024suffusion}
for coupled CFD--DEM suffusion under torsional shear, represents
the closest DEM analogue to the laboratory latex membrane. The
present formulation falls within the rigid-wall class and extends
it to undrained cyclic torsional loading. These citations now
appear in the DEM paragraph of Section~\ref{sec:introduction}.

\item \textbf{Specimen aspect ratio ($H=D_{\text{outer}}$).}
The Japanese Geotechnical Society standard for torsional shear on
hollow cylindrical specimens \citep{JGS0551_2020} does not prescribe
a specimen aspect ratio; it specifies equipment tolerances, test
procedure, and result processing in terms of generic $D_o$, $D_i$,
$H$ without constraining $H/D_{\text{outer}}$. In practice, both
the laboratory and DEM-HCA literature span
$H/D_{\text{outer}}\approx 1.0$ to $2.0$. In this range, the present work adopts
$H=D_{\text{outer}}=\SI{100}{mm}$, matching the laboratory-scale
HCA apparatus of \citet{Vargas2020} and the DEM-HCA
servo-formulation precursors \citet{Han2024jgs} and \citet{Ma2024}.
This choice is consistent with the scope of the present analysis,
which relies on specimen-averaged quantities (liquefaction
resistance in Fig.~\ref{fig:liq_resistance_curves} and
specimen-integrated fabric descriptors $Z_{m0}$ and $\alpha_0$ in
Section~\ref{sec:results}, all obtained by integration over the full
specimen boundary or the full contact population) rather than on
spatially resolved localisation features. A practical
consequence is that the particle count at fixed grading is roughly
halved compared with $H=2D_{\text{outer}}$, which is a material
saving given the parametric scope (Table~\ref{tab:case_matrix})
and the per-case cost of high-cycle undrained simulations near
liquefaction.
\end{enumerate}

\changed{Two manuscript locations updated:
(a) Section~\ref{sec:introduction} \manref{p.~2, l.~70--75}, DEM
paragraph: expanded to group the cited DEM-HCA precedents by radial
boundary treatment (rigid cylindrical walls \citep{Shi2021,Liu2021},
stacked-wall approximations of membrane compliance
\citep{LiGuoZhang2014,LiZhangGutierrez2015,LiChenGutierrez2017}, and
flexible bonded-sphere membranes
\citep{Song2024,Song2024suffusion}), and to position the present work
as an extension of the rigid-wall class to undrained cyclic loading.
(b) Section~\ref{subsec:dem_hca} \manref{p.~3, l.~119--121},
geometry sentence: revised to attribute the apparatus dimensions
explicitly to \citet{Vargas2020} and to note that the
$H = D_{\text{outer}}$ aspect ratio is also adopted by the DEM-HCA
servo-formulation precursors \citet{Han2024jgs} and \citet{Ma2024}.
The full $H/D_{\text{outer}}$-range survey and material-saving
rationale above are kept in this letter, since they expand on a
design-rationale question orthogonal to the paper's focus on $K_0$
effects and fabric evolution.}

\begin{revtext}
\textit{(Introduction, DEM paragraph):}
[\ldots] \textbf{DEM replication of HCA tests has been developed
under drained and monotonic loading, addressing principal stress
rotation, shear-band formation, and stress-path effects. Existing
DEM-HCA models can be grouped by the radial boundary treatment:
rigid cylindrical walls \citep{Shi2021,Liu2021}, stacked-wall
approximations of membrane compliance
\citep{LiGuoZhang2014,LiZhangGutierrez2015,LiChenGutierrez2017},
and, more recently, fully flexible bonded-sphere membranes
\citep{Song2024,Song2024suffusion}. However, extension to undrained
cyclic loading under HCA-type boundary conditions remains limited.}
\medskip

\textit{(Section~\ref{subsec:dem_hca}, geometry sentence):}
The specimen geometry (Fig.~\ref{fig:specimen_generation})
\textbf{follows the laboratory-scale HCA apparatus of
\citet{Vargas2020}, with inner diameter \SI{6}{cm}, outer diameter
\SI{10}{cm}, and specimen height \SI{10}{cm} ($H$ equal to the
outer diameter, also adopted by the DEM-HCA servo-formulation
precursors \citet{Han2024jgs} and \citet{Ma2024})}.
\end{revtext}

%% ---------------- R2.5 ----------------
\subsection*{R2.5 --- Stability of the servo near singular stiffness matrix}

\begin{rcomment}
The proposed servo mechanism may face numerical instability when the
stiffness matrix becomes near-singular due to the drastic loss of particle
contacts. How does the algorithm prevent numerical oscillations during
matrix inversion?
\end{rcomment}

\response
We thank the reviewer for raising this important point. In an
undrained DEM
specimen approaching liquefaction, the contact network thins
drastically and the aggregate wall stiffnesses collapse;
$\hat{\mathbf{K}}$'s entries shrink accordingly,
$\det\hat{\mathbf{K}}\!\to\!0$, and the servo coefficients
(the entries of $\hat{\mathbf{K}}^{-1}$) blow up. The failure
modes are division by zero at the matrix-inversion step and
divergence of the commanded wall velocities. This scenario arises
in two classes of undrained loading: cyclic loading as
$p^{\prime}\!\to\!0$ (cyclic liquefaction, at any density), and
monotonic loading on loose, contractive sand that can reach a
static / flow-liquefaction state along the same
$p^{\prime}\!\to\!0$ trajectory. Any servo that enters these
regimes therefore needs an explicit lower bound on the wall
stiffnesses to keep $\hat{\mathbf{K}}$ non-singular.

Our safeguard is a per-wall aggregate-stiffness clip that, by
construction, keeps every entry of $\hat{\mathbf{K}}$ bounded away
from zero, so division by zero and near-singular blow-up are
prevented at the level of matrix assembly rather than caught after
the fact. Concretely, each wall group's aggregate contact stiffness
(inner cylinder $K_{r,i}$, outer cylinder $K_{r,o}$, end platens
$K_z$) is floored at $10\%$ of its initial (end-of-consolidation)
value: once the measured value drops below that threshold, it is
clamped at the floor before $\hat{\mathbf{K}}$ is assembled. Both
diagonal entries of $\hat{\mathbf{K}}$ are sums of these per-wall
stiffnesses divided by wall areas; e.g.,
\begin{equation}
\hat{K}_{11} = -\frac{K_{r,i}}{A_i} - \frac{r}{R}\,\frac{K_{r,o}}{A_o},
\label{eq:rl_Khat11}
\end{equation}
with $A_i$, $A_o$ the corresponding wall areas. The clip therefore
keeps $\hat{K}_{11}$ (and analogously $\hat{K}_{22}$) strictly
negative and $\det\hat{\mathbf{K}}$ strictly non-zero regardless of
how many contacts are lost. An analogous $20\%$-of-initial clip on
the torsional moment of inertia protects the independent torsional
servo. Because $\hat{\mathbf{K}}$ is at most $2\times 2$, its
inverse is computed analytically in closed form at every timestep,
with no iterative linear solve involved.

\begin{figure}[htbp]
  \centering
  \includegraphics[width=0.5\linewidth]{figs/fig_r2_5.png}
  \caption{Aggregate wall stiffnesses $K_{r,i}$, $K_{r,o}$, $K_z$ and
    torsional moment of inertia $I_{rot}$ normalised by their
    end-of-consolidation values, throughout the same demanding
    cyclic run reported in R1.5
    ($D_r\approx 75\%$, $\Kzero=1.0$, CSR$=0.400$,
    constant-total-axial-stress boundary), plotted against the
    normalised cycle ratio $N_c/N_L$ ($\ru=0.95$ criterion). Black
    dashed: $10\%$ floor on each per-wall aggregate stiffness
    ($K_{r,i}$, $K_{r,o}$, $K_z$). Black dotted: $20\%$ floor on
    $I_{rot}$ in the independent torsional servo. Red dashed:
    liquefaction onset.}
  \label{fig:rl_floors}
\end{figure}

\begin{enumerate}[leftmargin=*,label=(\arabic*),topsep=4pt,itemsep=6pt]
\item \textbf{Observed clip behaviour and servo stability.}
The aggregate stiffnesses and moment of inertia for our most
demanding case ($D_r\approx 75\%$, $\Kzero=1.0$, CSR$=0.400$,
constant-total-axial-stress) are shown in
Fig.~\ref{fig:rl_floors}. Two observations follow:
\begin{enumerate}[leftmargin=2em,label=(\alph*),topsep=2pt,itemsep=3pt]
  \item The three per-wall stiffness ratios
    $K_{r,i}/K_{r,i}^{(0)}$, $K_{r,o}/K_{r,o}^{(0)}$, and
    $K_z/K_z^{(0)}$ each remain well above the $10\%$ floor
    throughout the run, with minimum values of $0.57$, $0.59$,
    and $0.47$, respectively. The clip on the radial--axial
    stiffness matrix $\hat{\mathbf{K}}$ is therefore a designed
    safety margin that is \emph{not} engaged in this run; the
    closed-form $2\times 2$ inversion remains well-conditioned
    on its own physical contact-network basis.
  \item The torsional moment of inertia $I_{rot}$ is more
    sensitive to the loss of blade--particle contacts as
    $\sigma_z^{\prime}\!\to\!0$. It decays cycle by cycle and
    reaches the $20\%$ floor at the deepest excursions of the
    post-liquefaction cycles, where the floor briefly engages
    and bounds the torsional servo coefficient
    $S_T=1/I_{rot}$.
\end{enumerate}
The servo runs stably through both regimes (cyclic
pre-liquefaction in the scalar-$\hat{\mathbf{K}}$ mode under
$dH/dt=0$, exercised in the $K_0$ suite that forms the paper's
core, and the full $2\times 2$ mode under constant-total-axial-stress,
exercised in the supplemental cyclic validation run reported here);
no high-frequency oscillation appears on any servo channel, and
the tracking-error ($e_r$, $e_z$) and specimen-volume traces
through these regimes are documented in the R2.6 validation
below.
\end{enumerate}

\changed{Section~\ref{subsec:cyclic_loading}
\manref{p.~8--9, l.~245--252}: a short sentence has been added after
the servo-coefficient derivation noting that the aggregate contact
stiffness on each wall is floored at $10\%$ of its initial
(end-of-consolidation) value before $\hat{\mathbf{K}}$ is assembled,
which keeps $\det\hat{\mathbf{K}}$ strictly non-zero and prevents
any near-singular behaviour in the closed-form inversion (an
analogous $20\%$ floor is applied to the torsional moment of inertia
in the independent torsional servo).}

\begin{revtext}
\textbf{In undrained loading, the contact network thins as the
effective confining stress decays, and the aggregate wall
stiffnesses $K_{r,i}$, $K_{r,o}$ and the end-platen stiffness $K_z$
that enter $\hat{\mathbf{K}}$ can drop towards zero, driving
$\det(\hat{\mathbf{K}})$ towards zero and the servo coefficients
towards arbitrarily large values. To prevent near-singular behaviour
in the closed-form inversion, each per-wall aggregate contact
stiffness is floored at $10\%$ of its end-of-consolidation value
before $\hat{\mathbf{K}}$ is assembled, which keeps every entry of
$\hat{\mathbf{K}}$ bounded away from zero by construction. An
analogous $20\%$ floor is applied to the rotational stiffness
$I_{rot}$ in the independent torsional servo
(\cref{eq:irot_scs}).} [\ldots] \textbf{Both floors function as
safety margins on the closed-form inversion and are exercised
explicitly in Section~\ref{subsec:numerical_verification}.}
\end{revtext}

%% ---------------- R2.6 ----------------
\subsection*{R2.6 --- Evolution of $e_r$, $e_z$, or volume during validation}

\begin{rcomment}
For the validation part, please add some figures to display the evolution
of $e_r$ and $e_z$ over time or the change in volume.
\end{rcomment}

\response
We thank the reviewer for this suggestion. The revised manuscript
includes a new methodology subsection
(Section~\ref{subsec:numerical_verification}) reporting the time
histories of the specimen-volume drift $V/V_0-1$ and the radial and
axial servo errors $e_r$, $e_z$ across a representative cyclic run
(Fig.~\ref{fig:rl_verification}). The verification case is
deliberately the most demanding in the simulation program
($D_r\approx 75\%$, $\Kzero=1.0$, CSR$=0.400$), and is run under the
constant-total-axial-stress axial boundary so that the full
$2\times 2$ coupled servo is exercised. The same verification case
has been added as a footnoted entry in Table~\ref{tab:case_matrix}.

\begin{figure}[htbp]
  \centering
  \includegraphics[width=\linewidth]{figs/fig_r2_6.pdf}
  \caption{Servo verification on the most demanding cyclic run
    ($D_r \approx 75\%$, $\Kzero = 1.0$, CSR $= 0.400$,
    constant-total-axial-stress boundary):
    (a)~specimen-volume drift $V/V_0 - 1$ confirming exact enforcement
    of the constant-volume constraint
    (Eq.~(\ref{eq:const_volume_constraint}));
    (b)~tracking errors of the radial and axial pressure-difference
    targets ($e_r$, $e_z$). Vertical dashed line: liquefaction onset
    ($N_c/N_L = 1$, $\ru = 0.95$). Panels~(a) and~(b) of
    Fig.~\ref{fig:numerical_verification} of the revised manuscript;
    panel~(c) of that figure (inertial number) is reproduced
    separately in R1.5 as Fig.~\ref{fig:rl_inertial}.}
  \label{fig:rl_verification}
\end{figure}

\begin{enumerate}[leftmargin=*,label=(\arabic*),topsep=4pt,itemsep=6pt]
\item \textbf{Volume conservation.}
Because the constant-volume constraint
(Eq.~(\ref{eq:const_volume_constraint})) is enforced kinematically
through the analytical relation between $dr/dt$, $dR/dt$, and $dH/dt$
rather than through a feedback loop, the specimen volume is
conserved to floating-point round-off: the relative drift
$V(t)/V_0-1$ stays below $10^{-7}$ throughout the run
(Fig.~\ref{fig:rl_verification}a), several orders of magnitude
smaller than the laboratory drift attributable to membrane
compliance and water compressibility.

\item \textbf{Servo tracking errors.}
The tracking errors $e_r=p_{dif,r}^{\prime}-p_{dif,r}^{tar}$ and
$e_z=p_{dif,z}^{\prime}-p_{dif,z}^{tar}$ are shown in
Fig.~\ref{fig:rl_verification}b. Both grow modestly as the effective
confining stress decays towards liquefaction but remain bounded
throughout, with peak magnitudes of only a few percent of the
initial mean effective stress. The closed-form $2\times 2$ inversion
of $\hat{\mathbf{K}}$ therefore continues to deliver
well-conditioned wall velocities even when the underlying contact
network is severely depleted, and the $10\%$ stiffness floor and
$20\%$ moment-of-inertia floor introduced in response to R2.5
keep $\det\hat{\mathbf{K}}$ bounded away from zero through the
post-liquefaction excursions.
\end{enumerate}

\changed{New Section~\ref{subsec:numerical_verification}
\manref{p.~10--11, l.~282--311} (``Numerical verification of the coupled
servo'') added at the end of the methodology, with
Fig.~\ref{fig:numerical_verification} \manref{p.~11} reporting
(a)~the relative volume drift, (b)~the radial and axial
servo-error traces, and (c)~the inertial number (panel~(c)
addressing R1.5). The verification case is listed as a footnoted
entry in Table~\ref{tab:case_matrix} \manref{p.~10}.}

\begin{revtext}
\textit{(New paragraph in
Section~\ref{subsec:numerical_verification}, ``Volume
conservation''):} \textbf{The constant-volume constraint is enforced
kinematically through the analytical relation between $dr/dt$,
$dR/dt$, and $dH/dt$ (\cref{eq:const_volume_constraint}), so the
specimen volume is expected to remain exactly conserved up to
floating-point round-off. Fig.~\ref{fig:numerical_verification}a
confirms this: the relative volume drift $V(t)/V_0 - 1$ stays below
$10^{-7}$ throughout the run, more than two orders of magnitude
smaller than the residual compressibility of laboratory pore
water.}
\medskip

\textit{(New paragraph, ``Boundary tracking errors''):}
\textbf{Fig.~\ref{fig:numerical_verification}b shows the tracking
errors for the radial and axial pressure-difference targets
($e_r$, $e_z$): both grow modestly as the effective confining stress
decays towards liquefaction (vertical dashed line), but remain
bounded throughout, with peak magnitudes a few percent of the
initial mean effective stress.} The closed-form $2 \times 2$
inversion of $\hat{\mathbf{K}}$ therefore continues to deliver
well-conditioned wall velocities even when the underlying contact
network is severely depleted [\ldots]
\end{revtext}

%% ---------------- R2.7 ----------------
\subsection*{R2.7 --- Contact-force anisotropy (3D rose diagrams)}

\begin{rcomment}
Fig.~14: You analysed contact density anisotropy, but why didn't you analyse
anisotropy for contact forces such as normal contact force? Please refer to
the papers mentioned before which plotted 3D rose diagrams of contact force
anisotropies.
\end{rcomment}

\response
We thank the reviewer for this important suggestion. We agree, and
have added contact-force anisotropy in two
complementary ways.

\begin{enumerate}[leftmargin=*,label=(\arabic*),topsep=4pt,itemsep=6pt]
\item \textbf{Eight-panel normal-force rose.}
A normal-force rose has been added as a new
Fig.~\ref{fig:contact_force_sequence} (immediately after the
count-weighted contact-density rose, now
Fig.~\ref{fig:contact_density_sequence}). For each of the eight
panels we extracted every ball-ball contact from
the PFC snapshot, projected the contact normal into the local
cylindrical frame ($r$, $\theta$, $z$), and binned the unit normals
on a $36\times18$ sphere with antipodal symmetry. Two fabric tensors
are then formed: the count-weighted
$\Phi_{ij}=(1/N_c)\sum_c n_i^c n_j^c$ and the normal-force-weighted
$\Phi^f_{ij}=(\sum_c f_n^c\,n_i^c n_j^c)/\!\sum_c f_n^c$
\citep{Zhou2022,Song2024}, the latter introduced as new
Eq.~\eqref{eq:fabric_tensor_force} in
Section~\ref{subsec:micro_fabric_results}. Each panel is annotated
with the deviator invariant $A$ and tilt $B$
(Eqs.~\eqref{eq:fabric_invariants} and~\eqref{eq:fabric_tilt}) and
the specimen-mean $\langle F_n\rangle$, so the count and force roses
can be compared on equal footing.

The figure shows that even when the count fabric is nearly isotropic
($A_c\approx 0.016$ for \Kzero{} = 2.0), the force fabric carries a much
stronger directional signature ($A_n=0.12$, $B_n\approx 90^\circ$) —
a small set of force-bearing contacts dominates load transfer.
Through liquefaction, $\langle F_n\rangle$ decays from
$\approx 0.57$\,N to $\approx 0.14$\,N (a $\sim 75\%$ reduction) while
$B_n$ in both specimens converges to $43$--$45^\circ$, aligning with
the rotating principal-stress direction. This corroborates the visual
observation in Fig.~\ref{fig:force_disp_sequence} (force chains) on a
quantitative basis.

\begin{figure}[htbp]
  \centering
  \includegraphics[width=\linewidth]{figs/fig_r2_7.pdf}
  \caption{Reproduced from Fig.~\ref{fig:contact_force_sequence} of
    the revised manuscript:
    normal-force rose for \Kzero{} = 0.5 (top row) and 2.0 (bottom row)
    at four stages of the liquefaction process ($N/N_L = 0.00$, 0.61,
    1.05, 1.07; CSR = 0.200). Each panel is annotated with the
    force-fabric deviator invariant $A_n$, the principal-axis tilt
    $B_n$ from the axial direction, and the specimen-mean
    $\langle F_n\rangle$, providing a per-panel quantitative
    counterpart to the original visual rose.}
  \label{fig:r2_7}
\end{figure}

\item \textbf{Mechanism plot over the full expanded dataset.}
The same count- and force-weighted decomposition is then applied to
the entire expanded dataset (five \Kzero{} states $\times$ two
relative densities, CSR = 0.300) at the consolidation end-state.
Figure~\ref{fig:fabric_mechanism}, now upgraded to a dual-axis plot,
overlays the count torsion-resisting fraction $1-\Phi_{rr}$ (left
axis) and the force-weighted counterpart $1-\Phi^f_{rr}$ (right
axis), with the underlying values listed in the new
Table~\ref{tab:fabric_diagonal} (Appendix~\ref{sec:appendix_fabric}).
Across $\Kzero=0.5\to 2.0$ at $D_r=90\%$, $1-\Phi_{rr}$ falls from
0.667 to 0.657 ($\sim 1.5\%$), whereas $1-\Phi^f_{rr}$ falls from
0.733 to 0.642 ($\sim 14\%$): a $\sim 9$-fold sensitivity gap. The
two density groups trace nearly identical curves for both
descriptors. This identifies the load-bearing subnetwork — rather
than the topological contact distribution alone — as the dominant
microscopic carrier of the \Kzero{}-dependent liquefaction
resistance, and now anchors both the mechanism discussion in
Section~\ref{subsec:micro_fabric_results} and the implications
paragraph in Section~\ref{subsec:limitations}.
\end{enumerate}

\changed{New Fig.~\ref{fig:contact_force_sequence} \manref{p.~20} and
dual-axis upgrade to Fig.~\ref{fig:fabric_mechanism} \manref{p.~23};
new Eq.~\eqref{eq:fabric_tensor_force} \manref{p.~12, l.~349--351} for
$\Phi^f_{ij}$; Table~\ref{tab:fabric_diagonal} \manref{p.~27} added
in Appendix~\ref{sec:appendix_fabric}; expanded paragraphs in
Section~\ref{subsec:micro_fabric_results}
\manref{p.~18--19, l.~494--503; p.~21, l.~548--560} and
Section~\ref{subsec:limitations} \manref{p.~21, l.~569--573};
\citet{Zhou2022} added to the bibliography.}

\begin{revtext}
\textit{(New equation in Section~\ref{subsec:micro_descriptors}):}
\textbf{$\Phi^f_{ij} = \dfrac{\sum_c f_n^c\, n_i^c\, n_j^c}{\sum_c f_n^c}$
weights each contact by its normal-force magnitude $f_n^c$ and
therefore reflects the directional partitioning of the load-bearing
network rather than the contact topology alone.}
\medskip

\textit{(New paragraph at
Fig.~\ref{fig:contact_force_sequence}):}
\textbf{Figure~\ref{fig:contact_force_sequence} weights the same
contact set by the normal contact force \citep{Zhou2022,Song2024}.
At the initial state, $A_c \approx 0.022$ for \Kzero{} = 0.5 and
$A_c \approx 0.016$ for \Kzero{} = 2.0, while the force-fabric
counterparts are several times larger ($A_n = 0.17$ and $0.12$, with
$B_n \approx 1^\circ$ and $90^\circ$ respectively).} [\ldots]
\textbf{Cyclic shearing reduces $\langle F_n \rangle$ from
$\approx 0.57$\,N to $\approx 0.14$\,N (a $\sim 75\%$ decay), while
$B_n$ converges to $43$--$45^\circ$ and $A_n$ rises to $\approx
0.18$--$0.20$ for both specimens} [\ldots]
\medskip

\textit{(Mechanism paragraph,
Section~\ref{subsec:micro_fabric_results}):} [\ldots]
\textbf{as \Kzero{} increases from 0.5 to 2.0 at $D_r = 90\%$,
$1-\Phi_{rr}$ falls from 0.667 to 0.657 (${\sim}1.5\%$), whereas
$1-\Phi^f_{rr}$ falls from 0.733 to 0.642 (${\sim}14\%$); the two
density groups trace nearly identical curves for both descriptors.}
[\ldots] \textbf{The much stronger \Kzero{}-sensitivity of $\Phi^f$
identifies the load-bearing network---rather than the contact
topology alone---as the dominant microscopic carrier of the
\Kzero{}-dependent liquefaction-resistance signal} [\ldots]
\end{revtext}

%% ---------------- R2.8 ----------------
\subsection*{R2.8 --- Quantifying the differences in Fig.~15 (now Fig.~\ref{fig:force_disp_sequence})}

\begin{rcomment}
Fig.~15: It is difficult to distinguish the difference in displacement
between $\Kzero=0.5$ and $2.0$ specimens when $N/N_L=0.00$ and $N/N_L=0.61$,
and it is also difficult to distinguish the difference in force chain when
$N/N_L=1.05$ and $N/N_L=1.07$. Please quantify the differences.
\end{rcomment}

\response
We thank the reviewer for flagging these specific instances. Both
sub-questions are now answered quantitatively through the new
normal-force rose (Fig.~\ref{fig:contact_force_sequence}, introduced
in response to R2.7), each of whose eight panels is annotated with
the force-fabric deviator
$A_n$, principal-axis tilt $B_n$, and specimen-mean
$\langle F_n\rangle$. The figure replaces the visual-only comparison
of the original Fig.~15 with a per-panel tagged comparison that
covers both reviewer-flagged regimes.

\begin{enumerate}[leftmargin=*,label=(\arabic*),topsep=4pt,itemsep=6pt]
\item \textbf{Force chain at $N/N_L = 1.05$ and $1.07$.}
At $N/N_L = 1.05$ the two specimens are essentially indistinguishable:
$A_n = 0.18$ vs.\ $0.20$, $B_n = 45^\circ$ vs.\ $43^\circ$,
$\langle F_n\rangle = 0.16$\,N vs.\ $0.17$\,N. At $N/N_L = 1.07$ the
convergence persists ($A_n = 0.20$ vs.\ $0.18$, $B_n = 43^\circ$ vs.\
$45^\circ$, $\langle F_n\rangle = 0.14$\,N for both). The reverse-
direction ordering between consecutive snapshots is the signature of
the post-liquefaction oscillation: each loading reversal swaps the
tilt direction of the residual force network. The visual observation
in Fig.~\ref{fig:force_disp_sequence} — that the force chains have
lost the imprint of the initial stress state once liquefaction is
reached — is now backed by these numerical values.

\item \textbf{Displacement at $N/N_L = 0.00$ and $0.61$.}
At these pre-liquefaction stages, particle displacements are
inherently small in both specimens because the contact network is
still load-supporting; visual similarity in the displacement panels
of Fig.~\ref{fig:force_disp_sequence} is therefore expected and is
itself part of the finding, not a sign of missing discrimination. The quantitative
discriminator at the same instants is not displacement but the
force-fabric: at $N/N_L = 0.00$ the \Kzero{} = 0.5 specimen carries
load axially ($A_n = 0.17$, $B_n \approx 1^\circ$,
$\langle F_n\rangle = 0.57$\,N) while the \Kzero{} = 2.0 specimen
carries load equatorially ($A_n = 0.12$, $B_n \approx 90^\circ$,
$\langle F_n\rangle = 0.57$\,N) — the two principal directions are
$\sim\!89^\circ$ apart even though the count-fabric and displacement
panels look nearly identical. By $N/N_L = 0.61$ the tilt has begun to
rotate toward the eventual $\sim\!45^\circ$ collapse direction
($B_n = 15^\circ$ for \Kzero{} = 0.5 and $B_n = 83^\circ$ for
\Kzero{} = 2.0), with $\langle F_n\rangle$ already reduced to
$\sim\!0.43$--$0.45$\,N. The new rose therefore distinguishes the
early-stage states quantitatively where the displacement field cannot,
and the displacement panels of Fig.~\ref{fig:force_disp_sequence} are
retained for their qualitative role of showing the kinematic
divergence that follows from the force-fabric divergence.
\end{enumerate}

\changed{R2.7 force-rose (Fig.~\ref{fig:contact_force_sequence}
\manref{p.~20}) cited as the quantitative discriminator at all four
time instants; per-panel $(A_n, B_n, \langle F_n\rangle)$ annotations
now present in the figure; displacement panels of
Fig.~\ref{fig:force_disp_sequence} \manref{p.~21} retained for
their kinematic-context role; explanatory paragraph extended in
Section~\ref{subsec:micro_fabric_results}
\manref{p.~19, l.~504--519}.}

\begin{revtext}
[\ldots] By $N/N_L = 0.61$ the displacements remain
indistinguishable, reflecting the small accumulated strain and high
shear stiffness at these pre-liquefaction stages; \textbf{the
quantitative discriminator between the two states is instead the
force-fabric tilt (Fig.~\ref{fig:contact_force_sequence}: $B_n =
1^\circ$ vs.\ $90^\circ$ at $N/N_L = 0.00$; $15^\circ$ vs.\
$83^\circ$ at 0.61), while the force chains thin out and darken
(the low-force end of the colorbar) and reorient toward the
rotating principal-stress direction.} The onset of liquefaction
($N/N_L = 1.05$) brings an abrupt change [\ldots] \textbf{Between
$N/N_L = 1.05$ and 1.07 the loading reverses and both visualisations
alternate accordingly: the displacement field flips to the opposite
direction, and the dominant force chains swing from one tilt to its
mirror in the $z$--$\theta$ plane ($B_n$ stays between $45^\circ$
and $43^\circ$ because, by definition, it is an unsigned tilt
magnitude). This reversal signature is shared by both \Kzero{}
states rather than discriminating between them.} [\ldots]
\end{revtext}

\clearpage

\section*{Additional minor clarifications}

The following minor edits were made proactively (not in response to a
specific reviewer comment) to remove residual ambiguities identified
during the revision pass. They are listed here so that the latexdiff
reflects them as intentional improvements rather than scope creep.

\begin{enumerate}[label=M\arabic*., leftmargin=*]
\item \textbf{Servo targets in laboratory total-stress terms
      (Section~\ref{subsec:cyclic_loading}).} The where-clauses of
      \cref{eq:cond1_radial,eq:cond2_axial}
      \manref{p.~6, l.~197--199; p.~6--7, l.~203--205} now state
      explicitly $p_{dif,r}^{tar}=p_o-p_i$ and
      $p_{dif,z}^{tar}=p_z-\sigma_r$, making the link from the formal
      definitions to the laboratory total-stress program clearer. The
      narrative paragraph above already established that the
      pore-pressure $u$ drops out of these differences; the equality
      is now reflected in the definitions themselves.

\item \textbf{Explicit servo state during AC
      (Section~\ref{subsec:preparation}).} The description of the
      anisotropic-consolidation stress path \manref{p.~4, l.~159--160}
      now states ``with $p_o^{\prime}=p_i^{\prime}$ held throughout
      (so $p_{dif,r}^{tar}=0$).'' This is symmetric with the
      cyclic-shear program opening
      ($p_o=p_i=\sigma_{r0}$, $p_{dif,r}^{tar}=0$) and removes any
      residual ambiguity about the radial servo state during AC.

\item \textbf{Table~\ref{tab:fabric_diagonal} caption
      (Appendix~\ref{sec:appendix_fabric}).} The earlier caption
      label ``(CSR = 0.300)'' was misleading because the tabulated
      values are evaluated at the consolidation end-state and are
      shared by all CSR cases at that $(D_r,\Kzero)$. The label
      \manref{p.~27} has been replaced with ``after anisotropic
      consolidation.'' The figure subcaptions in
      Section~\ref{subsec:micro_fabric_results} that legitimately
      refer to CSR$=0.300$ results (cyclic-stage observations) are
      unchanged.

\item \textbf{Flowchart annotations
      (Fig.~\ref{fig:servo_flowchart}).} Steps~4 and~5
      \manref{p.~9} now carry an italic sub-line indicating the
      constant-height reduction
      ($S_{rr}$ alone, \cref{eq:Srr_constantH_main}; $dr/dt = S_{rr}\,
      e_r/\Delta t$). This makes it visible from the figure alone
      that the parametric \Kzero{} study (constant height) uses the
      degenerate single-DOF form, while the full $2\times 2$ coupling
      is exercised in monotonic calibration and in the
      Section~\ref{subsec:numerical_verification} stress-control
      verification.

\item \textbf{Increment notation unified
      (\cref{eq:coupled_stiffness} and Appendix~\ref{sec:appendix_servo}).}
      The variational symbol $\delta p$ originally used on the LHS
      of the coupled stiffness relation \manref{p.~8, l.~234--235} has
      been replaced with the differential $dp$ (and likewise
      $\delta\sigma_r^{\prime}\to d\sigma_r^{\prime}$). Both sides of
      the tangent-stiffness relation now use the differential $d$
      uniformly, matching the $dr$, $dH$ on the RHS; the equation
      describes a real increment-to-increment mapping, not a
      virtual-work form.

\end{enumerate}

\clearpage

\section*{Closing}

We hope the revised manuscript now meets the journal's standards and the
reviewers' expectations. We remain grateful for the reviewers' careful
reading and valuable suggestions.

\bibliographystyle{plainnat}
\bibliography{refs}

\end{document}
